\chapter{Diabetes Tests}
\section*{What Is This Test?} Diabetes tests measure glucose regulation to screen for, diagnose, or monitor diabetes and prediabetes.
\section*{Alternative Names} Blood glucose testing; A1C; oral glucose tolerance test; diabetes screening.
\section*{Principle} Fasting glucose, random glucose, oral glucose tolerance, and hemoglobin A1C reflect current or average glycemia over different time periods.
\section*{Why This Test Is Ordered} Testing is used for symptoms, risk-factor screening, pregnancy, diagnosis, and monitoring treatment.
\section*{Primary vs Secondary Test} Glucose and A1C are primary diagnostic measures; urine glucose, ketones, autoantibodies, C-peptide, and kidney tests answer additional questions.
\section*{How This Test Is Performed} Blood is collected fasting, randomly, after a glucose drink, or for A1C; home meters monitor but do not replace laboratory diagnosis.
\section*{What Preparation Is Needed} Fasting and medication instructions vary. Follow the laboratory's directions and report anemia, hemoglobin variants, pregnancy, and recent illness.
\section*{Understanding Results} Results are interpreted using laboratory and guideline thresholds; acute illness and altered red-cell survival can affect results.
\section*{Follow-up Tests} Follow-up may include repeat diagnostic testing, kidney function, urine albumin, lipids, ketones, autoantibodies, or eye and foot assessment.
\section*{References} \begin{enumerate}\item MedlinePlus. Diabetes Tests.\item American Diabetes Association. Standards of Care in Diabetes.\end{enumerate}
\clearpage\section*{Understanding Results and Follow-up}
The laboratory report should identify the specimen, method, units, reference interval or decision limit, and whether the result is positive, negative, abnormal, or inconclusive. Reference intervals vary with age, sex, pregnancy, specimen type, assay, and laboratory; a result outside a range is not automatically a diagnosis, and a result within a range does not always exclude disease. Discuss unexpected results with the ordering clinician, who can decide whether repeat or confirmatory testing, treatment, or referral is appropriate. Do not change medicines or delay urgent care based on an isolated result.
\section*{Regulatory Status and Authorization Date}This chapter describes a test category, not a specific commercial product. FDA, CDSCO, EU IVDR, and MHRA approval or registration dates therefore cannot be assigned without the assay manufacturer, product name, instrument, and intended use. For laboratory-developed tests, an individual agency approval date may not exist. See the Preface for the interpretation of regulatory status.
\clearpage
